\documentclass[11pt]{article}
\usepackage[margin=1in]{geometry}
\usepackage{hyperref}
\usepackage{enumitem}
\usepackage{booktabs}
\usepackage{xcolor}
\title{ProbOS --- Month 3, Week 12\\Automotive / EV Safety Positioning}
\author{Nisong Monyimba}
\date{\today}

\begin{document}
\maketitle

\section*{Week 12 goal}
Produce accurate, well-researched positioning materials mapping
\texttt{BatteryModel2Cell}'s already-validated capabilities to
ISO 26262's automotive functional-safety framework.

\textbf{Critical distinction from Weeks 9--11, stated up front:} this
is NOT an engineering week. No new model, no new code, and no new
validation dataset are required -- the underlying battery model is
already Validated (Kim 2007, Sobol $S_1=0.457$ for Ea\_SEI). The work
here is entirely research and positioning: correctly understanding
ISO 26262's actual terminology and requirements (ASIL levels, Hazard
Analysis and Risk Assessment / HARA, safety case documentation) well
enough to map our existing capabilities onto them accurately, without
overclaiming compliance or certification we have not actually
obtained. Exit criteria are ``materials exist and are accurate,'' not
``tests pass.''

\section*{Day-by-day plan}

\subsection*{Day 1 -- ISO 26262 research (before writing any positioning material)}
\begin{itemize}[leftmargin=*]
  \item Research ISO 26262's actual structure: ASIL classification,
    the Hazard Analysis and Risk Assessment (HARA) process, and what
    documentation a safety case actually requires -- from real,
    citable sources, not assumed from general familiarity
  \item Identify precisely which parts of \texttt{BatteryModel2Cell}'s
    existing output (Sobol sensitivity indices, provenance-traced
    tail-risk outcomes, Monte Carlo percentile trajectories) map
    genuinely onto HARA-style hazard identification -- and be
    explicit about anything that does NOT map cleanly
  \item Confirm: are we claiming ProbOS \emph{could support} an
    ISO 26262 safety case, or that it \emph{has been used in} one?
    Only the former is currently true and must be the only claim made
\end{itemize}

\subsection*{Day 2 -- Draft the automotive/EV positioning brief}
\begin{itemize}[leftmargin=*]
  \item A dedicated technical brief for an automotive/EV safety
    engineering audience, explaining what \texttt{BatteryModel2Cell}
    already demonstrates and how it maps to ISO 26262 concepts
  \item Explicitly reuse the REAL, already-computed results (Kim 2007
    validation, $S_1=0.457$ for Ea\_SEI, the P95-worst-case-cell
    result) -- no new numbers are invented for this audience
\end{itemize}

\subsection*{Day 3 -- Outreach materials}
\begin{itemize}[leftmargin=*]
  \item Draft an outreach email template for automotive/EV safety
    engineers, analogous to \texttt{docs/enterprise\_pilot\_email.md}
    but addressed to this specific audience
  \item Keep the ask concrete and modest (a conversation to confirm
    or correct the assumed pain point), not an overreaching pitch
\end{itemize}

\subsection*{Day 4 -- Cross-check against the honest open questions}
\begin{itemize}[leftmargin=*]
  \item Re-read \texttt{docs/vision/main.tex}'s stated open questions
    and confirm the Week 12 materials do not implicitly resolve or
    paper over either of them
\end{itemize}

\subsection*{Day 5 -- Review pass}
\begin{itemize}[leftmargin=*]
  \item Full read-through checking for any claim not directly
    traceable to either an actual computed result already in the
    codebase, or an accurately cited external source
\end{itemize}

\subsection*{Day 6 -- Confirm no regressions}
\begin{itemize}[leftmargin=*]
  \item Run \texttt{pre\_week\_audit.sh} regardless, since this is a
    docs-only week -- confirms nothing was accidentally broken
\end{itemize}

\subsection*{Day 7 -- Documentation + retrospective}
\begin{itemize}[leftmargin=*]
  \item Update \texttt{docs/vision/main.tex}'s automotive section to
    note positioning materials are complete
  \item Retrospective appended to this document once complete
\end{itemize}

\section*{Exit criteria}
Accurate, well-researched automotive/EV positioning materials exist,
citing only real ISO 26262 concepts and real, already-computed ProbOS
results -- with no overclaimed certification, compliance, or
practitioner validation that has not actually occurred.


\section*{What was actually accomplished (retrospective)}

\begin{itemize}[leftmargin=*]
  \item Real ISO 26262 research conducted before drafting any
    positioning material: HARA (Part 3, Clause 7) confirmed to be a
    qualitative, structured engineering-judgment process (Severity
    S0--S3, Exposure E0--E4, Controllability, deriving ASIL and Safety
    Goals) -- NOT itself a quantitative simulation or UQ exercise.
  \item An honest, load-bearing distinction was established and
    maintained throughout both documents: ProbOS does not perform
    HARA, does not assign ASIL levels, and has not been used in an
    actual ISO 26262 submission.
  \item \texttt{docs/automotive\_ev\_positioning.md} and
    \texttt{docs/automotive\_ev\_outreach\_email.md} written, citing
    only real, already-computed ProbOS results.
  \item A verification script written to catch overclaiming language
    produced false alarms -- checker bugs (naive word matching
    without adequate negation-lookback distance; a substring check
    broken by a markdown line wrap; a vision-doc text replacement
    broken by an unexpected LaTeX line wrap), not real document
    problems. Each was diagnosed via direct grep before concluding
    the document itself was correct, rather than editing content to
    satisfy a possibly-broken check.
  \item Cross-checked against \texttt{docs/vision/main.tex}'s stated
    open questions: neither document claims a resolved competitive
    differentiation, nor implies practitioner validation has occurred.
\end{itemize}

\section*{What went right}
\begin{enumerate}[leftmargin=*]
  \item Real research surfaced the correct, load-bearing distinction
    between HARA-as-judgment-process and ProbOS-as-supporting-evidence.
  \item When automated verification flagged issues (three separate
    times), the actual document text was checked directly before
    concluding anything was wrong -- correctly distinguishing checker
    bugs from real content problems each time.
  \item This week correctly required no code changes.
\end{enumerate}

\section*{Numbers at Week 12 close}
\begin{itemize}[leftmargin=*]
  \item Python tests: 398/398 passing (unchanged, docs-only week)
  \item New documents: 2 (positioning brief, outreach email)
\end{itemize}

\section*{Carried into future weeks}
The outreach itself has not yet happened -- this week produced the
materials, not the practitioner validation itself.
\end{document}
